\subsubsection{Family-Based Analysis}
\begin{enumerate}
    \item \textbf{Calculation of family Centroids}
          
          The subroutine computes the centroid (mean expression vector) for each
          gene family in a gene expression dataset. It iterates over all
          families, selects the relevant genes for each family and calculates the
          mean vector for each family. There are two possible modes to use. To use
          all genes for calculation, select \texttt{all} mode. If you want to
          specify relevant genes for calculation, select \texttt{orthologs} mode.
          It is important that you provide a logical \texttt{ortholog\_set}
          array when using \texttt{orthologs} mode. Call the \texttt{group\_centroid}
          subroutine with the following arguments:
          
          \textbf{Input arguments:}
          \begin{itemize}
            \item \texttt{expression\_vectors real(real64)}: Expression vector matrix
                  with size \texttt{n\_axes $\times$ n\_genes}
                  
            \item \texttt{n\_axes integer(int32)}: Number of axes (tissues /
                  dimensions)
                  
            \item \texttt{n\_genes integer(int32)}: Total number of genes
                  
            \item \texttt{gene\_to\_family integer(int32)}: Gene-to-family mapping
                  array with the length of \texttt{n\_genes} (1-based indexing)
                  
            \item \texttt{n\_families integer(int32)}: Total number of gene families
                  
            \item \texttt{mode integer(int32)}: Calculation mode ("GROUP\_ALL" or "GROUP\_ORTHOLOGS") where "GROUP\_ALL" = 1 and "GROUP\_ORTHOLOGS" = 0 are parameters
                  
            \item (Optional) \texttt{ortholog\_set logical}: Logical array indicating
                  if a gene is part of the calculation subset (only used in \texttt{GROUP\_ORTHOLOGS}
                  mode)
          \end{itemize}
                  
          \textbf{Output arguments:}
          \begin{itemize}
            \item \texttt{centroid\_matrix real(real64)}: Matrix of calculated family
                  centroid vectors with size \texttt{n\_axes $\times$ n\_families}
                  
            \item \texttt{selected\_indices integer(int32)}: Array of gene indices
                  used for the calculation with length of \texttt{n\_genes}
                  
            \item \texttt{ierr integer(int32)}: Error code (for error details
                  please check \hyperref[sec:error-handling]{error handling section}
                  )
          \end{itemize}
          
          \begin{lstlisting}[language=Fortran, caption={Calculation of Family Centroids}]
call group_centroid(expression_vectors, n_axes, n_genes, &
            gene_to_family, n_families, centroid_matrix, &
            mode, selected_indices, ierr, ortholog_set)
          \end{lstlisting}
          
    \item \textbf{Compute Shift Vector Field}

        Computes shift vectors by subtracting corresponding family centroids from gene expression vectors. The output stores both centroids and shift vectors in a single matrix.
        
        \textbf{Mathematical Operation:}
        \[ \text{shift\_vectors}(d+i, j) = \text{expression\_vectors}(i, j) - \text{family\_centroids}(i, \text{gene\_to\_centroid}(j)) \]
        \[ \text{shift\_vectors}(1:d, j) = \text{family\_centroids}(:, \text{gene\_to\_centroid}(j)) \]
        
        \textbf{Arguments:}
        \begin{itemize}
            \item \texttt{d integer(int32)}: Expression vector dimension
            \item \texttt{n\_genes integer(int32)}: Total number of genes
            \item \texttt{n\_families integer(int32)}: Total number of families
            \item \texttt{expression\_vectors real(real64)}: Gene expression matrix [d × n\_genes]
            \item \texttt{family\_centroids real(real64)}: Family centroid matrix [d × n\_families]
            \item \texttt{gene\_to\_centroid integer(int32)}: Mapping from genes to family centroids [n\_genes]
            \item \texttt{shift\_vectors real(real64)}: Output matrix [2*d × n\_genes] storing centroids (rows 1-d) and shift vectors (rows d+1-2d)
            \item \texttt{ierr integer(int32)}: Error code
        \end{itemize}
        \begin{lstlisting}[language=Fortran, caption={Shift Vector Field Computation}]
call compute_shift_vector_field(d, n_genes, n_families, &
                    expression_vectors, family_centroids, &
                    gene_to_centroid, shift_vectors, ierr)
        \end{lstlisting}


        \textbf{Output Structure:}
        \begin{itemize}
            \item \texttt{shift\_vectors(1:d, :)}: Contains the family centroids for each gene
            \item \texttt{shift\_vectors(d+1:2*d, :)}: Contains the shift vectors (expression vector - family centroid)
        \end{itemize}
          
    \item \textbf{Calculate Mean Vector}
          
          Alternatively you can calculate a single mean vector for a given set of
          vectors by calling the \texttt{mean\_vector} subroutine directly with
          the following arguments:
          
          \textbf{Input arguments:}
          \begin{itemize}
            \item \texttt{expression\_vectors real(real64)}: Expression vector matrix
                  with size \texttt{n\_axes $\times$ n\_genes}
                  
            \item \texttt{n\_axes integer(int32)}: Number of axes (tissues /
                  dimensions)
                  
            \item \texttt{n\_genes integer(int32)}: Total number of genes
                  
            \item \texttt{gene\_indices integer(int32)}: An array containing the
                  column indices of the selected genes in \texttt{expression\_vectors}
                  
            \item \texttt{n\_selected\_genes integer(int32)}: Total number of genes
                  in the current family to be calculated
          \end{itemize}
          
          \textbf{Output arguments:}
          \begin{itemize}
            \item \texttt{centroid real(real64):} Calculated mean vector (centroid)
                  with length of \texttt{n\_axes}
                  
            \item \texttt{ierr integer(int32)}: Error code (for error details
                  please check \hyperref[sec:error-handling]{error handling section}
                  )
          \end{itemize}
          \begin{lstlisting}[language=Fortran, caption={Mean Vector Calculation}]
call mean_vector(expression_vectors, n_axes, n_genes, gene_indices,&
         n_selected_genes, centroid, ierr)
          \end{lstlisting}
          
    \item \textbf{Calculate Distance to Centroid}
          
          For each gene in the expression vector matrix, the subroutine computes
          the Euclidean distance to its corresponding family centroid and
          returns a distances array containing the distance value for each
          expression vector to its family centroid. Call the \texttt{distance\_to\_centroid}
          subroutine with the following arguments:
          
          \textbf{Input arguments:}
          \begin{itemize}
            \item \texttt{n\_genes integer(int32)}: Total number of genes
                  
            \item \texttt{n\_families integer(int32)}: Total number of gene families
                  
            \item \texttt{genes real(real64)}: Gene expression matrix (\texttt{d
                  $\times$ n\_genes})
            
            \item \texttt{centroids real(real64)}: Family centroid matrix (\texttt{d
                  $\times$ n\_genes})
            
            \item \texttt{gene\_to\_fam integer(int32)}: Gene-to-family mapping array
                  with the length of \texttt{n\_genes} (1-based indexing)
                  
            \item \texttt{d integer(int32)}: Expression vector dimension
          \end{itemize}
          \textbf{Output arguments:}
          \begin{itemize}
            \item \texttt{distances real(real64)}: Array of distances for each gene
                  with length of \texttt{n\_genes}
          \end{itemize}
          \begin{lstlisting}[language=Fortran, caption={Distance to Centroid Calculation}]
call distance_to_centroid(n_genes, n_families, genes, centroids, &
                    gene_to_fam, distances, d)
          \end{lstlisting}
          
          \begin{enumerate}
            \item \textbf{Calculate Euclidean Distance}
                          
                  Alternatively you can calculate a single Euclidean distance between
                  two vectors manually by calling \texttt{euclidean\_distance}
                  subroutine with the following arguments:
                  
                  \textbf{Input arguments:}
                  \begin{itemize}
                    \item \texttt{vec1 real(real64):} First expression vector as
                          array with length of dimension \texttt{d}
                          
                    \item \texttt{vec2 real(real64):} Second expression vector as
                          array with length of dimension \texttt{d}
                          
                    \item \texttt{d integer(int32):} Dimension value
                  \end{itemize}
                  
                  \textbf{Output arguments:}
                  \begin{itemize}
                    \item \texttt{result real(real64):} Euclidean distance value
                          between \texttt{vec1} and \texttt{vec2}
                  \end{itemize}
                  \begin{lstlisting}[language=Fortran, caption={Euclidean distance calculation}]
call euclidean_distance(vec1, vec2, d, result)
                  \end{lstlisting}
          \end{enumerate}
          
    \item \textbf{Detect Gene Outliers}
          
          The subroutine first computes family scaling, then calculates relative
          distance indices (RDI) and identifies outliers based on the specified
          percentile threshold. It returns a boolean array indicating whether each
          gene is an outlier. Call the \texttt{detect\_outliers} subroutine with
          the following arguments:
          
          \textbf{Input arguments:}
          \begin{itemize}
            \item \texttt{n\_genes integer(int32)}: Total number of genes
                  
            \item \texttt{n\_families integer(int32)}: Total number of gene families
                  
            \item \texttt{distances real(real64)}: Array of distances for each gene
                  to its centroid with length of \texttt{n\_genes}
                  
            \item \texttt{gene\_to\_fam integer(int32)}: Gene-to-family mapping array
                  with the length of \texttt{n\_genes} (1-based indexing)
                  
            \item \textbf{(Optional)} \texttt{percentile real(real64)}:
                  Percentile threshold value for outlier detection (default is 95)
          \end{itemize}
          
          \textbf{Input / Output arguments:}
          \begin{itemize}
            \item \texttt{work\_array real(real64)}: Work array for sorting with
                  length of \texttt{n\_genes}
                  
            \item \texttt{perm integer(int32)}: Permutation array for sorting with
                  length of \texttt{n\_genes}
                  
            \item \texttt{stack\_left integer(int32)}: Stack array for left indices
                  during sorting with length of \texttt{n\_genes}
                  
            \item \texttt{stack\_right integer(int32)}: Stack array for right indices
                  during sorting with length of \texttt{n\_genes}
                  
            \item \texttt{loess\_x real(real64)}: Reference x-coordinates for loess
                  smoothing with length of \texttt{n\_families}
                  
            \item \texttt{loess\_y real(real64)}: Reference y-coordinates for loess
                  smoothing with length of \texttt{n\_families}
                  
            \item \texttt{loess\_n integer(int32)}: Indices of reference points used
                  for smoothing with length of \texttt{n\_families}
          \end{itemize}
          
          \textbf{Output arguments:}
          \begin{itemize}
            \item \texttt{is\_outlier logical}: Boolean array with length of \texttt{n\_genes}
                  indicating whether each gene is an outlier
                  
            \item \texttt{ierr integer(int32)}: Error code (for error details
                  please check \hyperref[sec:error-handling]{error handling section}
                  )
          \end{itemize}
          
          \begin{lstlisting}[language=Fortran, caption={Gene Outlier Calculation}]
call detect_outliers(n_genes, n_families, distances, gene_to_fam, &
            work_array, perm, stack_left, stack_right, &
            is_outlier, loess_x, loess_y, loess_n, ierr, &
            percentile)
          \end{lstlisting}
          
          \textbf{Alternatively you can calculate each step of \texttt{detect\_outliers}
          manually with the following subroutine:}
          \begin{enumerate}
            \item \textbf{Compute Family Scaling}
                  
                  The subroutine calculates a scaling factor for each gene family by
                  computing the median and standard deviation of gene-to-centroid distances
                  within each family, then applies LOESS smoothing to these values. Call
                  the \texttt{compute\_family\_scaling} subroutine with the
                  following arguments:
                  
                  \textbf{Input arguments:}
                  \begin{itemize}
                    \item \texttt{n\_genes integer(int32)}: Total number of genes
                          
                    \item \texttt{n\_families integer(int32)}: Total number of gene families
                          
                    \item \texttt{distances real(real64)}: Array of distances for each
                          gene to its centroid with length of \texttt{n\_genes}
                          
                    \item \texttt{gene\_to\_fam integer(int32)}: Gene-to-family mapping
                          array with the length of \texttt{n\_genes} (1-based indexing)
                  \end{itemize}
                  
                  \textbf{Input / Output arguments:}
                  \begin{itemize}
                    \item \texttt{loess\_x real(real64)}: Reference x-coordinates for
                          loess smoothing with length of \texttt{n\_families}
                          
                    \item \texttt{loess\_y real(real64)}: Reference y-coordinates for
                          loess smoothing with length of \texttt{n\_families}
                          
                    \item \texttt{indices\_used integer(int32)}: Indices of reference
                          points used for smoothing with length of \texttt{n\_families}
                          
                    \item \texttt{perm\_tmp integer(int32)}: Permutation array for sorting
                          with length of \texttt{n\_genes}
                  \end{itemize}
                  
                  \textbf{Output arguments:}
                  \begin{itemize}
                    \item \texttt{dscale real(real64)}: Array of scaling factors per
                          family with length of \texttt{n\_families}
                          
                    \item \texttt{family\_distances real(real64)}: Temporary work array
                          with length of \texttt{n\_genes} to store distances of genes
                          within each family during calculation
                          
                    \item \texttt{ierr integer(int32)}: Error code (for error
                          details please check
                          \hyperref[sec:error-handling]{error handling section} )
                  \end{itemize}
                  \begin{lstlisting}[language=Fortran, caption={Compute Family Scaling}]
call compute_family_scaling(n_genes, n_families, distances, &
                gene_to_fam, dscale, loess_x, loess_y, &
                indices_used, perm_tmp, stack_left_tmp, &
                stack_right_tmp, family_distances, ierr)
                  \end{lstlisting}
                  
            \item \textbf{Compute Family Scaling Helper}
                  
                  This helper subroutine allocates internal arrays for easier usage and
                  calls \texttt{compute\_family\_scaling}. Call the \texttt{compute\_family\_scaling\_alloc}
                  subroutine with the following arguments:
                  
                  \textbf{Input arguments:}
                  \begin{itemize}
                    \item \texttt{n\_genes integer(int32)}: Total number of genes
                          
                    \item \texttt{n\_families integer(int32)}: Total number of gene families
                          
                    \item \texttt{distances real(real64)}: Array of distances for each
                          gene to its centroid with length of \texttt{n\_genes}
                          
                    \item \texttt{gene\_to\_fam integer(int32)}: Gene-to-family mapping
                          array with the length of \texttt{n\_genes} (1-based indexing)
                  \end{itemize}
                  
                  \textbf{Input / Output arguments:}
                  \begin{itemize}
                    \item \texttt{loess\_x real(real64)}: Reference x-coordinates for
                          loess smoothing with length of \texttt{n\_families}
                          
                    \item \texttt{loess\_y real(real64)}: Reference y-coordinates for
                          loess smoothing with length of \texttt{n\_families}
                          
                    \item \texttt{indices\_used integer(int32)}: Indices of reference
                          points used for smoothing with length of \texttt{n\_families}
                  \end{itemize}
                  
                  \textbf{Output arguments:}
                  \begin{itemize}
                    \item \texttt{dscale real(real64)}: Array of scaling factors per
                          family with length of \texttt{n\_families}
                          
                    \item \texttt{ierr integer(int32)}: Error code (for error
                          details please check
                          \hyperref[sec:error-handling]{error handling section} )
                  \end{itemize}
                  \begin{lstlisting}[language=Fortran, caption={Compute Family Scaling}]
call compute_family_scaling_alloc(n_genes, n_families, distances, &
                        gene_to_fam, dscale, loess_x, &
                        loess_y, indices_used, ierr)
                  \end{lstlisting}
                  
            \item \textbf{Compute Relative Distance Index (RDI)}
                  
                  The subroutine calculates the Relative Distance Index (RDI) for each
                  gene by dividing its Eculidean distance to the family centroid by the
                  scaling factor of its family. Call the \texttt{compute\_rdi} subroutine
                  with the following arguments:
                  
                  \textbf{Input arguments:}
                  \begin{itemize}
                    \item \texttt{n\_genes integer(int32)}: Total number of genes
                          
                    \item \texttt{distances real(real64)}: Array of distances for each
                          gene to its centroid with length of \texttt{n\_genes}
                          
                    \item \texttt{gene\_to\_fam integer(int32)}: Gene-to-family mapping
                          array with the length of \texttt{n\_genes} (1-based indexing)
                          
                    \item \texttt{dscale real(real64)}: Array of scaling factors per
                          family with length of \texttt{n\_families}
                  \end{itemize}
                  
                  \textbf{Input / Output arguments:}
                  \begin{itemize}
                    \item \texttt{sorted\_rdi real(real64)}: Work array for sorting with
                          length of \texttt{n\_genes}
                          
                    \item \texttt{perm integer(int32)}: Pre-initialized (1: \texttt{n\_genes})
                          permutation array for sorting with length of \texttt{n\_genes}
                          
                    \item \texttt{stack\_left integer(int32)}: Stack array for sorting
                          with length of \texttt{n\_genes}
                          
                    \item \texttt{stack\_right integer(int32)}: Stack array for sorting
                          with length of \texttt{n\_genes}
                  \end{itemize}
                  
                  \textbf{Output arguments:}
                  \begin{itemize}
                    \item \texttt{rdi real(real64)}: Array of RDI values for each gene
                          with length of \texttt{n\_genes}
                  \end{itemize}
                  \begin{lstlisting}[language=Fortran, caption={Compute Relative Distance Index (RDI)}]
call compute_rdi(n_genes, distances, gene_to_fam, dscale, rdi, &
        sorted_rdi, perm, stack_left, stack_right)
                  \end{lstlisting}
                  
            \item \textbf{Identify Outliers}
                  
                  The subroutine identifies outliers based on the top percentile of RDI
                  values. Call the \texttt{identify\_outliers} subroutine with the following
                  arguments:
                  
                  \textbf{Input arguments:}
                  \begin{itemize}
                    \item \texttt{n\_genes integer(int32)}: Total number of genes
                          
                    \item \texttt{rdi real(real64)}: Array of RDI values for each gene
                          with length of \texttt{n\_genes}
                          
                    \item \texttt{sorted\_rdi real(real64)}: Array of sorted RDI values
                          for each gene with legnth of \texttt{n\_genes} - \textbf{this
                            array must be filtered to remove negatives and be sorted in
                          ascending order!}
                          
                    \item \textbf{(Optional)} \texttt{percentile real(real64)}:
                          Percentile threshold value for outlier detection (default is 95)
                  \end{itemize}
                  
                  \textbf{Output arguments:}
                  \begin{itemize}
                    \item \texttt{is\_outlier logical}: Boolean array with length of
                          \texttt{n\_genes} indicating whether each gene is an outlier
                          
                    \item \texttt{threshold}: Actual threshold value used for outlier
                          detection
                  \end{itemize}
                  \begin{lstlisting}[language=Fortran, caption={Identify Outliers}]
call identify_outliers(n_genes, rdi, sorted_rdi, is_outlier, &
                threshold, percentile)
                  \end{lstlisting}
          \end{enumerate}
                  
    \item \textbf{Compute Normalized Tissue Versatility}
          
          The subroutine calculates a normalized tissue versatility score for selected
          gene expression vectors. It is based on the angle between each gene
          expression vector and the space diagonal. Versatility is normalized to
          [0, 1], where 0 means uniform expression 1 means expression in only one
          axis. Call the \texttt{compute\_tissue\_versatility} subroutine with the
          following parameters:
          
          \textbf{Input arguments:}
          \begin{itemize}
            \item \texttt{n\_genes integer(int32)}: Total number of genes
                  
            \item \texttt{n\_vectors integer(int32)}: Total number of expression
                  vectors
                  
            \item \texttt{n\_selected\_axes integer(int32)}: Total number of selected
                  axes
                  
            \item \texttt{n\_selected\_vectors integer(int32)}: Total number of selected
                  vectors
                  
            \item \texttt{expression\_vectors real(real64)}: Expression vector matrix
                  with size \texttt{n\_axes $\times$ n\_vectors}
                  
            \item \texttt{exp\_vecs\_selection\_index logical}: Logical array with
                  length of \texttt{n\_vectors} to specify which vectors should be
                  included in the calculation
                  
            \item \texttt{axes\_selection logical}: Logical array with length of
                  \texttt{n\_axes} to specify which axes should be included in the
                  calculation
          \end{itemize}
          
          \textbf{Output arguments:}
          \begin{itemize}
            \item \texttt{tissue\_versatilities real(real64)}: Array with the length
                  of \texttt{n\_selected\_vectors} containing the calculated tissue
                  versatilities
                  
            \item \texttt{tissue\_angles\_deg real(real64)}: Array with the length
                  of \texttt{n\_selected\_vectors} containing the calculated angles
                  in degrees
                  
            \item \texttt{ierr integer(int32)}: Error code (for error details
                  please check \hyperref[sec:error-handling]{error handling section}
                  )
          \end{itemize}
          \begin{lstlisting}[language=Fortran, caption={Tissue Versatility Calculation}]
call compute_tissue_versatility(n_axes, n_vectors, &
        expression_vectors, exp_vecs_selection_index, &
        n_selected_vectors, axes_selection, &
        n_selected_axes, tissue_versatilities, &
        tissue_angles_deg, ierr)
          \end{lstlisting}
\end{enumerate}
